\section{Experience-Dependent Structural Plasticity in Knowledge Graphs: The Bridge Twin Engine}

\textbf{Authors}: Bryan Alexander Buchorn  
\textbf{Affili\-ations}: Independent Researcher, Las Vegas, NV, USA  
\textbf{Status}: v1.2.0 (Hardened Enterprise)  
\textbf{Date}: March 2026

---

\subsubsection{Abstract}
We introduce the \textbf{Bridge Twin Engine}, a mechanism for autonomous topological evolution in Knowledge Graphs (KGs) based on reasoning experience. While traditional KGs maintain a static structure, the Bridge Twin Engine implements a graph-theoretic analog to biological Long-Term Potent\-iation (LTP) \cite{hebb1949, bipoo1998}. By mater\-ializing "Twin Nodes" to act as structural relays across frequently traversed community boundaries, the system short-circuits latent paths and optimizes its own geometry for future inference. We formalize the poten\-tiation rules based on usage frequency and semantic salience, and provide a pruning mechanism (LTD) to maintain topological sanity. The v1.2.0 release incor\-porates an atomic re-mapping protocol to synchronize relays with background community re-parti\-tioning, effectively solving the "Zombie Bridge" problem in dynamic envir\-onments.

\subsubsection{1. Introd\-uction}
The efficiency of multi-hop reasoning in graphs is funda\-mentally constrained by the diameter and sparsity of the network. In large-scale KGs, seman\-tically related concepts often reside in disparate community partitions, neces\-sitating high-latency "Bridge" traversals. We propose that a graph should not be a static artifact, but a plastic entity that reshapes its topology to favor successful reasoning chains.

\subsubsection{2. Methodology}

\paragraph{2.1 The Potent\-iation Rule}
A structural relay node $v'_{twin}$ is mater\-ialized for node $v$ in destination community $C_{dest}$ if:
\begin{enumerate}
\item \textbf{Frequency}: Traversal $(u \to v)$ occurs $\geq n_{min}$ times.
\item \textbf{Alignment}: $\cos(\vec{e}_v, \vec{c}_{dest}) \geq \sigma$, where $\vec{c}_{dest}$ is the community centroid.

\end{enumerate}
\paragraph{2.2 Materi\-alization and Reflection}
Upon poten\-tiation, the engine performs a "Reflection" operation:
\begin{itemize}
\item Assign $v'_{twin}$ to $C_{dest}$.
\item Establish a \texttt{BRIDGE\_TWIN} binding edge $(v, v'_{twin})$.
\item Mirror all of $v$'s edges that terminate in $C_{dest}$ onto $v'_{twin}$.

\end{itemize}
\subsubsection{3. Structural Maintenance: The LTD Analog}
To prevent the "Inform\-ational Sprawl" of mater\-ialized nodes, we implement a temporal decay function:
\begin{equation}
c_t = c_0 \cdot \lambda^{\Delta t}
\end{equation}
where $\lambda$ is the decay constant and $\Delta t$ is the time since last usage. Edges falling below a confidence threshold $c_{min}$ are pruned during the \textbf{REM Cycle} (SPEC\_007).

\subsubsection{4. Conclusion}
The Bridge Twin Engine provides a biolo\-gically-inspired framework for self-optimizing Knowledge Graphs. By trans\-forming reasoning history into physical graph structure, it enables sub-millisecond multi-hop reasoning on incre\-asingly complex network topologies.

---
\textbf{References}
\begin{enumerate}
\item Hebb, D. O. (1949). The organ\-ization of behavior: A neuro\-psychological theory.
\item Bi, G. Q., \& Poo, M. M. (1998). Synaptic modif\-ications in cultured hippocampal neurons. Journal of Neuros\-cience.
\item Markram, H., et al. (1997). Regulation of synaptic efficacy by predictions of spike timing. Science.
\item Newman, M. E. (2006). Modularity and community structure in networks. PNAS.
\item Abbott, L. F., \& Nelson, S. B. (2000). Synaptic plasticity: taming the beast. Nature Neuros\-cience.
\item Caporale, N., \& Dan, Y. (2008). Spike timing-dependent plasticity: a Hebbian learning rule. Annual Review of Neuros\-cience.
\item Buchorn, B. A., \& Sonnet, C. (2026). Bridge Twin Relays in CEREBRUM. SPEC\_003.md.

\end{enumerate}